% GENERATED FILE. Edit canonical lessons, then run npm run generate:books.
% canonical-source-hash: fnv1a64:877361d10913b4ad
% canonical-lessons: PT-C430-pingar, PT-C430-ralo, PT-C430-entupir, PT-C430-desentupir, PT-C430-autoclismo

\chapter{A casa de banho não funciona}
\label{ch:pt-a-casa-de-banho-nao-funciona}
\hlchaptermodality{430}

\begin{chapteropening}
\textbf{By the end of this chapter:} \emph{I can report problems in the bathroom: dripping, the drain, a blockage, getting it unblocked and the cistern.}

Complete the last lesson of chapter 430: I can report problems in the bathroom: dripping, the drain, a blockage, getting it unblocked and the cistern.
\end{chapteropening}

\section[pingar]{pingar — to drip}

\label{lesson:PT-C430-pingar}



\begin{quote}
Before the new one: say the Portuguese for the coat hanger, then the Portuguese for the bathrobe.
\end{quote}

\subsection*{You'll want to know: pingar}
\textbf{pingar} — \textquotedblleft{}to drip\textquotedblright{}.

\begin{quote}
A torneira pinga toda a noite.
\end{quote}

\begin{quote}
\emph{The tap drips all night.}
\end{quote}

\emph{Um pingo} — a drop. In Portugal \emph{um pingo} is also a small coffee with milk.

\begin{cousinweb}
The origin of \emph{pingar} is not settled; it may imitate the sound of a falling drop, and EtymDB has no entry for it.
\end{cousinweb}

\begin{grammarlens}[title={What you've built}]
One more word for telling, explaining and describing at B1.
\end{grammarlens}

\subsection*{Your turn}
\begin{itemize}
\raggedright
  \item \emph{Say it:} \emph{pingar}
  \item \emph{Say it:} \emph{pingar}, once more
  \item \emph{Say it:} a sentence of your own with \emph{pingar}
  \item \emph{Recall:} say the Portuguese for the coat hanger, then the Portuguese for the bathrobe, then say \emph{pingar} again
\end{itemize}

\begin{tcolorbox}[breakable,colback=teal!4,colframe=teal!35!black,title={Before you move on}]
What does pingar mean? (\textquotedblleft{}To drip\textquotedblright{}.) Say it once more.
\end{tcolorbox}
\section[o ralo]{o ralo — the drain (in a shower or sink)}

\label{lesson:PT-C430-ralo}



\begin{quote}
Before the new one: say the Portuguese for the bathrobe, then the Portuguese for to drip.
\end{quote}

\subsection*{You'll want to know: o ralo}
\textbf{o ralo} — \textquotedblleft{}the drain (in a shower or sink)\textquotedblright{}.

\begin{quote}
A água não desce pelo ralo do duche.
\end{quote}

\begin{quote}
\emph{The water does not go down the shower drain.}
\end{quote}

As an adjective \emph{ralo} means thin, sparse: \emph{cabelo ralo}.

\begin{cousinweb}
EtymDB traces \emph{ralo} to Latin \emph{rallum} 'scraper', but the origin is not settled; the adjective \emph{ralo} 'thin' comes from Latin \emph{rārus}.
\end{cousinweb}

\begin{grammarlens}[title={What you've built}]
One more word for telling, explaining and describing at B1.
\end{grammarlens}

\subsection*{Your turn}
\begin{itemize}
\raggedright
  \item \emph{Say it:} \emph{o ralo}
  \item \emph{Say it:} \emph{o ralo}, once more
  \item \emph{Say it:} a sentence of your own with \emph{o ralo}
  \item \emph{Recall:} say the Portuguese for the bathrobe, then the Portuguese for to drip, then say \emph{o ralo} again
\end{itemize}

\begin{tcolorbox}[breakable,colback=teal!4,colframe=teal!35!black,title={Before you move on}]
What does o ralo mean? (\textquotedblleft{}The drain (in a shower or sink)\textquotedblright{}.) Say it once more.
\end{tcolorbox}
\section[entupir]{entupir — to block, to clog}

\label{lesson:PT-C430-entupir}



\begin{quote}
Before the new one: say the Portuguese for to drip, then the Portuguese for the drain.
\end{quote}

\subsection*{You'll want to know: entupir}
\textbf{entupir} — \textquotedblleft{}to block, to clog\textquotedblright{}.

\begin{quote}
O lavatório entupiu outra vez.
\end{quote}

\begin{quote}
\emph{The washbasin got blocked again.}
\end{quote}

\emph{Estar entupido} — to be blocked; a blocked nose too.

\begin{cousinweb}
The origin of \emph{entupir} is not settled, and EtymDB has no entry for it.
\end{cousinweb}

\begin{grammarlens}[title={What you've built}]
One more word for telling, explaining and describing at B1.
\end{grammarlens}

\subsection*{Your turn}
\begin{itemize}
\raggedright
  \item \emph{Say it:} \emph{entupir}
  \item \emph{Say it:} \emph{entupir}, once more
  \item \emph{Say it:} a sentence of your own with \emph{entupir}
  \item \emph{Recall:} say the Portuguese for to drip, then the Portuguese for the drain, then say \emph{entupir} again
\end{itemize}

\begin{tcolorbox}[breakable,colback=teal!4,colframe=teal!35!black,title={Before you move on}]
What does entupir mean? (\textquotedblleft{}To block, to clog\textquotedblright{}.) Say it once more.
\end{tcolorbox}
\section[desentupir]{desentupir — to unblock, to unclog}

\label{lesson:PT-C430-desentupir}



\begin{quote}
Before the new one: say the Portuguese for the drain, then the Portuguese for to block.
\end{quote}

\subsection*{You'll want to know: desentupir}
\textbf{desentupir} — \textquotedblleft{}to unblock, to unclog\textquotedblright{}.

\begin{quote}
O canalizador desentupiu o ralo numa hora.
\end{quote}

\begin{quote}
\emph{The plumber unblocked the drain in an hour.}
\end{quote}

Built from \emph{des-} and \emph{entupir}.

\begin{cousinweb}
\emph{desentupir} is built from \emph{des-} 'un-' and \emph{entupir} 'to block', whose own origin is not settled; EtymDB has no entry for either word.
\end{cousinweb}

\begin{grammarlens}[title={What you've built}]
One more word for telling, explaining and describing at B1.
\end{grammarlens}

\subsection*{Your turn}
\begin{itemize}
\raggedright
  \item \emph{Say it:} \emph{desentupir}
  \item \emph{Say it:} \emph{desentupir}, once more
  \item \emph{Say it:} a sentence of your own with \emph{desentupir}
  \item \emph{Recall:} say the Portuguese for the drain, then the Portuguese for to block, then say \emph{desentupir} again
\end{itemize}

\begin{tcolorbox}[breakable,colback=teal!4,colframe=teal!35!black,title={Before you move on}]
What does desentupir mean? (\textquotedblleft{}To unblock, to unclog\textquotedblright{}.) Say it once more.
\end{tcolorbox}
\section[o autoclismo]{o autoclismo — the toilet cistern, the flush}

\label{lesson:PT-C430-autoclismo}



\begin{quote}
Before the new one: say the Portuguese for to block, then the Portuguese for to unblock.
\end{quote}

\subsection*{You'll want to know: o autoclismo}
\textbf{o autoclismo} — \textquotedblleft{}the toilet cistern, the flush\textquotedblright{}.

\begin{quote}
O autoclismo não para de correr.
\end{quote}

\begin{quote}
\emph{The cistern keeps running.}
\end{quote}

\emph{Puxar o autoclismo} — to flush. In Brazil: \emph{a descarga}.

\begin{cousinweb}
\emph{autoclismo} is a learned coinage from Greek \emph{autós} 'self' and \emph{klýsma} 'liquid for washing out' (EtymDB).
\end{cousinweb}

\begin{grammarlens}[title={What you've built}]
One more word for telling, explaining and describing at B1.
\end{grammarlens}

\subsection*{Your turn}
\begin{itemize}
\raggedright
  \item \emph{Say it:} \emph{o autoclismo}
  \item \emph{Say it:} \emph{o autoclismo}, once more
  \item \emph{Say it:} a sentence of your own with \emph{o autoclismo}
  \item \emph{Recall:} say the Portuguese for to block, then the Portuguese for to unblock, then say \emph{o autoclismo} again
\end{itemize}

\begin{tcolorbox}[breakable,colback=teal!4,colframe=teal!35!black,title={Before you move on}]
What does o autoclismo mean? (\textquotedblleft{}The toilet cistern, the flush\textquotedblright{}.) Say it once more.
\end{tcolorbox}
